\section{Light --- Lib --- Logs}
%%% generated by the blueprint pipeline from TCSlib.Programs.Light.Lib.Logs
\subsection*{Light --- Lib --- Logs: Introduction}

Declarations of the Lean module \texttt{TCSlib.Programs.Light.Lib.Logs}.

\subsection*{Light --- Lib --- Logs: Results}

\begin{definition}[Argument variable index]
\lean{Light.Logs.Arg}
\label{Light.Logs.Arg}
\leanok
The argument variable is assigned index $0$.
\end{definition}

\begin{definition}[Candidate variable index]
\lean{Light.Logs.Cand}
\label{Light.Logs.Cand}
\leanok
\uses{Light.Logs.Arg}
The candidate variable is assigned index $1$.
\end{definition}

\begin{definition}[Power variable index]
\lean{Light.Logs.Power}
\label{Light.Logs.Power}
\leanok
\uses{Light.Logs.Arg}
The power variable is assigned index $2$.
\end{definition}

\begin{definition}[Body of the logarithm routine]
\lean{Light.log2Body}
\label{Light.log2Body}
\leanok
\uses{Light.Logs.Arg, Light.Logs.Cand, Light.Logs.Power, Light.Stmt, Light.k, Light.v}
The statement sequence initializes a candidate variable to $0$ and a power variable to
$2$. While the power is at most the input variable, it doubles the power and increments
the candidate by $1$; it then assigns the candidate to the input variable.
\end{definition}

\begin{definition}[Logarithm time bound]
\lean{Light.log2Time}
\label{Light.log2Time}
\leanok
For every natural number $x$, define the time bound $T(x)=14\lfloor\log_2 x\rfloor+12$.
\end{definition}

\begin{theorem}[Logarithm routine meets its time bound]
\lean{Light.log2_meets}
\label{Light.log2_meets}
\leanok
\uses{Light.Ends.next, Light.Ends.whileBlock, Light.Meets, Light.Meets.of_body, Light.frame, Light.log2Body, Light.log2Time, Light.update_frame_setLocal}
\difficulty{5}
Let $\mathrm{lim}$ be machine limits, $P$ a program, $d,p,x\in\mathbb N$, and
$\mu:\mathbb N\to\mathbb Z$. Suppose procedure $p$ in $P$ is the routine that computes
the base-two logarithm, and the word limit is at least $2x+2$. Then calling procedure
$p$ at depth $d$ with the single argument $x$ and initial memory $\mu$ has an execution
within time $14\lfloor\log_2 x\rfloor+12$ whose returned value is $\lfloor\log_2
x\rfloor$ and whose final memory is $\mu$.
\end{theorem}

\begin{definition}[Candidate and power loop body]
\lean{Light.clog2Body}
\label{Light.clog2Body}
\leanok
\uses{Light.Logs.Arg, Light.Logs.Cand, Light.Logs.Power, Light.Stmt, Light.k, Light.v}
The body assigns $0$ to a candidate variable and $1$ to a power variable. While the
power variable is less than an argument variable, each iteration doubles the power
variable and increments the candidate variable by $1$; the body then assigns the
candidate variable’s value to the argument variable.
\end{definition}

\begin{definition}[Logarithmic time bound]
\lean{Light.clog2Time}
\label{Light.clog2Time}
\leanok
For each natural number $x$, define the time bound to be $12\lceil\log_2 x\rceil+10$,
where $\lceil\log_2 x\rceil$ is the least natural number $k$ such that $x\le 2^k$.
\end{definition}

\begin{theorem}[Logarithmic computation by a program]
\lean{Light.clog2_meets}
\label{Light.clog2_meets}
\leanok
\uses{Light.Ends.next, Light.Ends.whileBlock, Light.Meets, Light.Meets.of_body, Light.clog2Body, Light.clog2Time, Light.frame, Light.update_frame_setLocal}
\difficulty{4}
Let a program $P$ contain the specified logarithm-computation body at procedure index
$p$, let $\mu:\mathbb{N}\to\mathbb{Z}$ be an initial memory, and let $x\in\mathbb{N}$.
If $2x+1\leq\mathrm{word}(\mathrm{lim})$, then calling procedure $p$ with input $x$ from
the initial state whose location values are given by the one-element argument list $[x]$
and whose memory is $\mu$ can finish within $12\,\mathrm{clog}(2,x)+10$ steps, returning
$\mathrm{clog}(2,x)$ in location $0$ and leaving memory equal to $\mu$.
\end{theorem}

\begin{theorem}[A cube above every natural number]
\lean{Light.exists_le_cube}
\label{Light.exists_le_cube}
\leanok
\difficulty{1}
For every natural number $n$, there is a natural number $s$ such that $n\le s^3$.
\end{theorem}

\begin{definition}[Least cube upper bound]
\lean{Light.cbrtLeast}
\label{Light.cbrtLeast}
\leanok
\uses{Light.exists_le_cube}
For each natural number $n$, define $r$ to be the least natural number whose cube is at
least $n$.
\end{definition}

\begin{theorem}[Cube bound for the least cube root]
\lean{Light.le_cbrtLeast_pow}
\label{Light.le_cbrtLeast_pow}
\leanok
\uses{Light.cbrtLeast, Light.exists_le_cube}
\difficulty{1}
For every natural number $n$, if $r$ is the least natural number such that $n\le r^3$,
then $n\le r^3$.
\end{theorem}

\begin{theorem}[Cube bound below the least cube root]
\lean{Light.pow_lt_of_lt_cbrtLeast}
\label{Light.pow_lt_of_lt_cbrtLeast}
\leanok
\uses{Light.cbrtLeast, Light.exists_le_cube}
\difficulty{1}
For natural numbers $n$ and $s$, if $s$ is less than the least natural number whose cube
is at least $n$, then $s^3<n$.
\end{theorem}

\begin{theorem}[Least cube-root bound]
\lean{Light.cbrtLeast_le}
\label{Light.cbrtLeast_le}
\leanok
\uses{Light.cbrtLeast, Light.exists_le_cube}
\difficulty{1}
Let $n,s$ be natural numbers. If $n\le s^3$, then the least natural number $r$
satisfying $n\le r^3$ is at most $s$.
\end{theorem}

\begin{theorem}[Least cube root bounded by its input]
\lean{Light.cbrtLeast_le_self}
\label{Light.cbrtLeast_le_self}
\leanok
\uses{Light.cbrtLeast, Light.cbrtLeast_le}
\difficulty{1}
For every natural number $n$, the least natural number $r$ satisfying $n\le r^3$ is at
most $n$.
\end{theorem}

\begin{theorem}[Positive lower bound for the least cube root]
\lean{Light.one_le_cbrtLeast}
\label{Light.one_le_cbrtLeast}
\leanok
\uses{Light.cbrtLeast, Light.le_cbrtLeast_pow}
\difficulty{3}
For a natural number $n$ with $1\le n$, let $r$ be the least natural number such that
$n\le r^3$. Then $1\le r$.
\end{theorem}

\begin{theorem}[Cube of the least cube root]
\lean{Light.cbrtLeast_pow_le}
\label{Light.cbrtLeast_pow_le}
\leanok
\uses{Light.cbrtLeast, Light.pow_lt_of_lt_cbrtLeast}
\difficulty{4}
For every natural number $n$, if $r$ is the least natural number satisfying $n\le r^3$,
then $r^3\le 8n$.
\end{theorem}

\begin{theorem}[Bounds for squares and cubes]
\lean{Light.abs_cube_le}
\label{Light.abs_cube_le}
\leanok
\uses{Light.cbrtLeast, Light.cbrtLeast_pow_le}
\difficulty{5}
Let $n,i$ be natural numbers, and let $L$ be an integer word bound. If $8n+1\le L$ and
$i$ is at most the least natural number $r$ such that $n\le r^3$, then $|(i: \mathbb
Z)i|\le L$ and $|(i: \mathbb Z)i^2|\le L$.
\end{theorem}

\begin{definition}[Cube ceiling procedure]
\lean{Light.cbrtCeilBody}
\label{Light.cbrtCeilBody}
\leanok
\uses{Light.Logs.Arg, Light.Logs.Cand, Light.Stmt, Light.k, Light.v}
Define a program that sets a candidate to $0$, repeatedly increments it by $1$ while its
cube is less than the input, and then replaces the input with the candidate.
\end{definition}

\begin{definition}[Cube ceiling time bound]
\lean{Light.cbrtCeilTime}
\label{Light.cbrtCeilTime}
\leanok
\uses{Light.cbrtLeast}
For each natural number $n$, define the time bound to be $12r+12$, where $r$ is the
least natural number whose cube is at least $n$.
\end{definition}

\begin{theorem}[Execution of the ceiling cube-root program]
\lean{Light.cbrtCeil_meets}
\label{Light.cbrtCeil_meets}
\leanok
\uses{Light.Ends.next, Light.Ends.whileBlock, Light.Meets, Light.Meets.of_body, Light.abs_cube_le, Light.cbrtCeilBody, Light.cbrtCeilTime, Light.cbrtLeast, Light.cbrtLeast_le_self, Light.frame, Light.le_cbrtLeast_pow, Light.pow_lt_of_lt_cbrtLeast, Light.update_frame_setLocal}
\difficulty{6}
Let $P$ be a finite program, let $p,n,d$ be natural numbers, let $\mu:\mathbb
N\to\mathbb Z$, and let the limits have word bound $W$. Suppose the program at position
$p$ is the specified loop program that starts a candidate at $0$, increments it while
its cube is less than the input, and then writes the candidate to the argument location.
If $8n+1\le W$, then this program call, with input list $[n]$, memory $\mu$, and depth
parameter $d$, finishes within $12r+12$ steps, where $r$ is the least natural number
satisfying $n\le r^3$; it returns $r$ in location $0$ and leaves memory equal to $\mu$.
\end{theorem}

\begin{definition}[Halving procedure]
\lean{Light.halfBody}
\label{Light.halfBody}
\leanok
\uses{Light.Logs.Arg, Light.Logs.Cand, Light.Stmt, Light.k, Light.v}
Define a program that sets a candidate to $0$, repeatedly increments it by $1$ while
twice the candidate is less than the input, and then replaces the input with the
candidate.
\end{definition}

\begin{definition}[Halving time bound]
\lean{Light.halfTime}
\label{Light.halfTime}
\leanok
For each natural number $h$, define the time bound to be
$10\left\lfloor\frac{h+1}{2}\right\rfloor+10$.
\end{definition}

\begin{theorem}[Execution of the halving program]
\lean{Light.half_meets}
\label{Light.half_meets}
\leanok
\uses{Light.Ends.next, Light.Ends.whileBlock, Light.Meets, Light.Meets.of_body, Light.frame, Light.halfBody, Light.halfTime, Light.update_frame_setLocal}
\difficulty{5}
Let a program $P$ contain the halving program at position $p$, let $h$ be a natural
number, and let $\mu:\mathbb N\to\mathbb Z$ be a memory state. Suppose
$h+1\leq\mathrm{word}(\mathrm{lim})$. Starting at depth $d$ with input frame $[h]$ and
memory $\mu$, the program meets the time bound $10\lfloor(h+1)/2\rfloor+10$ at position
$p$: it has an execution whose output at location $0$ is the integer
$\lfloor(h+1)/2\rfloor$ and whose memory remains $\mu$.
\end{theorem}
